\section{Evaluation}\label{sec:eval}

We evaluate \toolname against the general LTL miner of Neider and Gavran~\cite{NeiderG18} on a suite of 20 GR(1) benchmarks.
All experiments run on a single core of an Apple M-series machine with a 10-minute timeout per benchmark.

\subsection{Benchmark Suite}

No standard GR(1) trace-mining benchmarks exist.
We construct a suite of 20 specifications covering five dimensions:

\begin{enumerate}
\item \textbf{Variable count} (2--6 atomic propositions).
\item \textbf{Propositional depth} of the components ($D \in \{1,2,3\}$).
\item \textbf{Template variant}: justice-only (J), with initial conditions (IJ), with safety (SJ), and full GR(1) (ISJ).
\item \textbf{Multiple justice/guarantee conditions} ($m,k \in \{1,2\}$).
\item \textbf{Propositional connectives}: variables only, negation, conjunction, disjunction, and their nesting.
\end{enumerate}

For each benchmark we fix a ground-truth GR(1) formula $\varphi^*$ and generate traces by random sampling: we draw 5\,000 random lasso traces of length 4--6 with a random lasso start, evaluate $\varphi^*$, and retain 15 positive and 10 negative traces.
Both tools receive identical trace sets.

\paragraph{Baseline configuration.}
The baseline searches over the full LTL operator set $\{\always, \eventually, \Next, \until, \land, \lor, \lnot, \rightarrow\}$ using the DAG encoding of~\cite{NeiderG18}, iterating DAG size $D = 1, 2, \ldots, 10$.
\toolname receives the GR(1) template (number of justice/guarantee conditions and whether init/safety slots are active) and iterates $D = 1, 2, \ldots$ over the shared propositional depth.

\subsection{Results}

Table~\ref{tab:results} reports, for each benchmark, the time to find a separating formula and the DAG depth at which a solution is found.

\begin{table}[t]
\centering
\caption{Evaluation results. Times in seconds; ``TO'' = 10\,min timeout.
$D_e$: expected propositional depth. Tmpl: template variant.
Speedup is the ratio of baseline to \toolname wall-clock time.}\label{tab:results}
\vspace{2pt}
\small
\setlength{\tabcolsep}{3pt}
\begin{tabular}{@{}l r r l r r r r r@{}}
\toprule
Benchmark & $n$ & $D_e$ & Tmpl & \multicolumn{2}{c}{\toolname} & \multicolumn{2}{c}{Baseline~\cite{NeiderG18}} & Speedup \\
\cmidrule(lr){5-6}\cmidrule(lr){7-8}
& & & & Time & $D$ & Time & $D$ & \\
\midrule
\multicolumn{9}{@{}l}{\emph{Justice-only, variable scaling}} \\
B01 & 2 & 1 & J   & 0.06 & 1 & 5.42 & 5 & 98$\times$ \\
B02 & 3 & 1 & J   & 0.03 & 1 & 48.15 & 6 & 1\,719$\times$ \\
B03 & 4 & 1 & J   & 0.03 & 1 & 0.60 & 3 & 18$\times$ \\
B04 & 5 & 1 & J   & 0.04 & 1 & 31.93 & 6 & 823$\times$ \\
B20 & 6 & 1 & J   & 0.04 & 1 & 6.32 & 5 & 148$\times$ \\
\addlinespace
\multicolumn{9}{@{}l}{\emph{Justice-only, propositional depth}} \\
B05 & 2 & 2 & J   & 0.11 & 2 & 0.74 & 3 & 7$\times$ \\
B06 & 4 & 2 & J   & 0.17 & 2 & 0.59 & 3 & 4$\times$ \\
B07 & 3 & 2 & J   & 0.43 & 3 & TO & -- & ${>}1\,398\times$ \\
B08 & 4 & 2 & J   & 0.15 & 1 & 6.63 & 5 & 44$\times$ \\
B09 & 3 & 3 & J   & 1.07 & 4 & TO & -- & ${>}562\times$ \\
\addlinespace
\multicolumn{9}{@{}l}{\emph{Multiple justice/guarantee conditions}} \\
B10 & 3 & 1 & J   & 0.04 & 1 & 4.92 & 5 & 121$\times$ \\
B11 & 3 & 1 & J   & 0.04 & 1 & TO & -- & ${>}15\,036\times$ \\
B12 & 4 & 1 & J   & 0.09 & 1 & TO & -- & ${>}6\,771\times$ \\
\addlinespace
\multicolumn{9}{@{}l}{\emph{Init, safety, and full template}} \\
B13 & 3 & 1 & IJ  & 0.05 & 1 & 108.86 & 7 & 2\,191$\times$ \\
B14 & 3 & 1 & SJ  & 0.05 & 1 & 0.58 & 3 & 11$\times$ \\
B15 & 4 & 1 & ISJ & 0.09 & 1 & 0.60 & 3 & 7$\times$ \\
B16 & 5 & 1 & ISJ & 0.11 & 1 & 0.59 & 3 & 6$\times$ \\
B17 & 4 & 1 & ISJ & 0.12 & 1 & 0.59 & 3 & 5$\times$ \\
B18 & 3 & 2 & ISJ & 0.07 & 1 & 0.57 & 3 & 8$\times$ \\
\addlinespace
\multicolumn{9}{@{}l}{\emph{Complex}} \\
B19 & 4 & 2 & J   & 0.65 & 3 & 346.73 & 7 & 535$\times$ \\
\bottomrule
\end{tabular}
\end{table}

\paragraph{Solving rate.}
\toolname solves all 20 benchmarks, while the baseline times out on 4 (B07, B09, B11, B12).
The timeouts occur on benchmarks with conjunctions in justice conditions (B07, B09) and multiple guarantee conditions (B11, B12)---precisely the cases where the GR(1) structure is hardest to recover from an unconstrained LTL search.

\paragraph{Speedup.}
On the 16 benchmarks both tools solve, \toolname is $3.5\times$--$2{,}191\times$ faster (geometric mean: $45\times$).
Including the 4~timeouts, lower-bound speedups exceed $15{,}000\times$.
The largest gains arise when the baseline must search deep DAGs ($D \geq 5$) to represent what \toolname captures at $D = 1$.

\paragraph{DAG depth.}
The baseline requires DAG depths 3--7, because it must construct the full temporal structure (e.g., $\always(\eventually x_1)$ needs $D \geq 3$).
\toolname finds most solutions at $D = 1$ since only the propositional content is searched and many specifications use atomic propositions directly.

\paragraph{Formula quality.}
Beyond speed, the GR(1) template ensures that every mined formula is in the realisable GR(1) fragment and directly usable for synthesis.
The baseline, unconstrained, frequently produces formulas outside any known synthesis fragment.
Among the 16 benchmarks the baseline solves, we observe three patterns:
\begin{itemize}
\item \emph{Meaningful GR(1)-like formulas} (6/16): e.g., $\always(x_0 \rightarrow \eventually x_1)$ for benchmark B01.  These capture a reactive relationship but are not in canonical GR(1) form.
\item \emph{Overly general formulas} (5/16): e.g., $\always\eventually x_1$, which satisfies the traces but ignores the assumption side entirely.
\item \emph{Spurious separators} (5/16): e.g., $x_1 \until \lnot x_1$, which is consistent with the traces but encodes no meaningful reactive specification.  This formula holds on any trace where $x_1$ eventually becomes false, which happens to separate the given traces by coincidence.
\end{itemize}
By contrast, every formula \toolname produces has the canonical GR(1) shape and correctly identifies both the assumption and guarantee components.
